%=====================================================================
% 第 6 节  重新得到几何 Langlands 对应
% 原文：Section 6. Recovering the geometric Langlands correspondence
%       （arXiv:1702.06499v2）
%=====================================================================
\section{重新得到几何 Langlands 对应}\label{Hecke1}

\begin{keytakeaways}
\begin{itemize}
  \item \textbf{物理图像：}带面算符的 $\CS$ 类理论在 Nekrasov--Shatashvili 极限 $\ep_2\ra0$ 下，
        其配分函数的指数按 $1/\ep_2$ 发散；发散系数 $\mathcal{Y}_\si$ 是 opers 的 Lagrangian
        子流形（在复 Fenchel--Nielsen 坐标下）的生成函数，而指数中的有限部分给出 Hitchin
        本征方程的解 $\Psi$。这正是几何 Langlands 对应在物理上的具体体现。
  \item \textbf{数学工具：}仿射李代数共形块在临界水平极限下的渐近展开、opers 与复
        Fenchel--Nielsen 坐标、Hitchin 纤维化与 Hitchin Hamiltonians，以及 Cecotti--Vafa 的
        $tt^*$ 方程（其 $tt^*$-联络是平坦的 $SL(2,\mathbb{C})$-联络）。
  \item \textbf{关键结论：}余四维面算符的加入等价于在 $t$ 处插入退化场；$\ep_2\ra0$ 时配分函数
        按 $\psi_v(t)$ 因子化，这正是 Hecke 本征值性质。$\psi_v(t)$ 的单值性定义一个 oper
        局部系 $\chi$，$tt^*$-联络与之规范等价。
  \item \textbf{技术要点：}$\ep_2\ra0$ 与 CFT 临界水平极限的对应由 $k=-h^{\vee}-\ep_2/\ep_1$
        实现；$tt^*$-联络要成为 oper 联络，须其在基 $\Pi_i^v$ 下的反全纯部分幂零，且由水平性
        条件给出的二阶方程在 $C_\tau$ 上无奇点。于是 Beilinson--Drinfeld 的 $\CD$-模由
        $D_r-E_r$ 生成的理想之商给出。
\end{itemize}
\end{keytakeaways}

\bigskip
\noindent\rule{\textwidth}{0.6pt}
\bigskip

本节的目标是详细阐述：在面算符存在的情形下，AGT 对应的各种推广形式与几何 Langlands
对应之间的关系。我们将论证，在 Nekrasov--Shatashvili 极限 $\ep_2\ra 0$ 下会出现一种直接的
关联。为此我们先指出：配分函数 $\CZ_\si(a;x,\tau;\ep_1,\ep_2)$ 的渐近行为中含有若干解，
它们满足描述几何 Langlands 对应中 $\CD$-模结构的 Hitchin 本征方程。然后我们讨论：
如何从带面算符的 $\CS$ 类理论的 Nekrasov--Shatashvili 极限这一角度，来理解相应
$\CD$-模的关键性质——Hecke 本征值性质。这些现象在以 $\CM_H(G)$ 为目标流形的 $2d$ sigma
模型中的可能解释，将在下一节讨论。

\subsection{Nekrasov--Shatashvili 极限}

极限 $\ep_2\ra 0$ 依 \cite{NS} 而称为 Nekrasov--Shatashvili 极限。关系式
$k=-h^{\vee}-\ep_2/\ep_1$ 把它等同于 CFT 中的临界水平极限。在 CFT 中已知 \cite{T10,FGT}：
仿射李代数共形块的手征配分函数在临界水平极限下表现为
\begin{equation}\label{NSlim}
\mathcal{Z}_\si(a;x,\tau;\ep_1,\ep_2)\sim
e^{-\frac{1}{\ep_2}\mathcal{Y}_\si(a;\tau;\ep_1)}
\Psi(a;x,\tau;\ep_1)\,\big(1+\CO(\ep_2)\big)\,,
\end{equation}
其中
\begin{itemize}
\item $\mathcal{Y}_\si(a;\tau;\ep_1)$
是 opers 的 Lagrangian 子流形 ${\rm Op}_{\fsl_2}(C_\tau)$ 的生成函数，该对象已在 \cite{NRS}
的一个相关语境中考虑过；而 $a$ 是 ${\rm Op}_{\fsl_2}(C_\tau)$ 的复 Fenchel--Nielsen
坐标，由沿 $C$ 上决定裤分解 $\si$ 的一组曲线的和乐之迹定义。
\item $\Psi(a;x,\tau;\ep_1)$ 是 Hitchin Hamiltonians 的本征函数，$D_r\Psi=E_r\Psi$，其本征值
$E_r=E_r(a;\tau;\ep_1)$ 由 $\mathcal{Y}_\si(a;\tau;\ep_1)$ 按
$E_r=\partial_{\tau_r}\mathcal{Y}_\si(a;\tau;\ep_1)$ 给出。
\end{itemize}
式\eqref{NSlim} 的行为可看作几何 Langlands 对应的一种具体体现：特征标簇中复
Fenchel--Nielsen 坐标为 $(a,\pa_a\mathcal{Y}_\si(a;\tau;\ep_1))$ 的点，可以表示为
$\ep_1$-opers $(\CE_{o},\nabla'_{\ep_1,u})=
 \big(\CE_{o},\big(\ep_1\pa_z+ \big(
 \begin{smallmatrix} 0 & u \\ 1 & 0\end{smallmatrix}\big)\big)dz\big)$ 的和乐。
对一个满足 $u=u^{\sharp}+\sum_{r=1}^{3g-3+n}E_r\vartheta_r(x)$ 的 oper（这里 $u^\sharp$ 定义某个
固定的参考 oper，而 $\vartheta_r(x)(dz)^2$，$r=1,\dots,3g-3+n$ 构成 $H^0(C,K^2)$ 的一组基），
我们可以联系起方程组 $D_r\Psi=E_r\Psi$，$r=1,\dots,3g-3+n$，它定义了几何 Langlands
对应中与 oper $(\CE_{o},\nabla'_{\ep_1,u})$ 相联系的 $\CD$-模。

形如式\eqref{NSlim} 的行为，从带面算符的 $\CS$ 类理论的角度是容易理解的。
在极限 $\ep_2\ra 0$ 下，人们会预期：沿一个在 $\ep_2\ra 0$ 时仍保持有效紧致的子流形
$\CM^2$ 修改规范场的边界条件，不会影响该极限下的领头奇异行为。

\subsection{CFT 中的 Hecke 本征值性质}\label{Hecke-CFT}

当两类面算符同时存在时，可以定义配分函数的进一步推广
$\CZ_v(a;x,\tau;t;\ep_1,\ep_2)$：除此前考虑过的全部变量之外，它额外依赖于一个变量 $t$，
而 $t$ 可以等同于 $C$ 上一点的坐标，即一条“原初”（primordial）M2-膜的端点。上文已论证：
这类配分函数在 CFT 中表示为共形块的手征配分函数，并在 $t$ 处多出一个退化场的插入。
额外的标号 $v$ 指涉某个基的元素，该基张成给出 $4d+2d$ 拉格朗日描述的 GLSM 的真空空间。
在固定一种裤分解的选取之后，我们暂时略去标号 $\si$。

由 CFT \cite{T10,FGT} 可知，配分函数 $\CZ_v(a;x,\tau;t;\ep_1,\ep_2)$ 在极限 $\ep_2\ra 0$ 下表现为
\begin{equation}\label{NSlim+}
\CZ(a;x,\tau;t;\ep_1,\ep_2)\sim
e^{-\frac{1}{\ep_2}\mathcal{Y}(a;\tau;\ep_1)}
\Psi(x;a,\tau;\ep_1)\,\psi_v(t;a,\tau;\ep_1)\big(1+\CO(\ep_2)\big)\,,
\end{equation}
其中 $\mathcal{Y}(a;\tau;\ep_1)$ 与 $\Psi(x;a,\tau;\ep_1)$ 就是式\eqref{NSlim} 中出现的那两个
函数，而 $\psi_v(t)=\psi_v(t;a,\tau;\ep_1)$ 张成微分方程
$(\ep_1^2\partial_t^2-u(t))\psi_v(t)=0$ 之解空间的一组基，这里 $u=u(t)$ 把该 oper 联络表示为
$\ep_1\partial_t+\big(\begin{smallmatrix} 0& u \\ 1 & 0\end{smallmatrix}\big)$。
指标 $v=+,-$ 标记解空间的一组基。

$\psi_v(t)$ 之所以能作为一个因子被分离出来，与 Hecke 本征值性质有关，这一点已在
\cite{T10} 中指出。与前面一样，我们可以把 $\mathcal{Z}_v^*$ 解释为重叠
$\langle \de_x^{\sst (2)},\xi_a^{}\rangle$，其中 $\xi_a\in\CH^{\sst (2)}_a$。若式\eqref{NSlim+}
对所有 $x$ 成立，则对一切 $\xi_a\in\CH^{\sst (2)}_a$ 都有类似的因子化。此外，因子
$\psi_v(t)$ 并不依赖于为把配分函数表为 $\langle \de_x^{\sst (2)},\xi_a^{}\rangle$ 所选用的
向量 $\xi_a\in\CH^{\sst (2)}_a$。把 $\psi_v(t)$ 解析延拓所生成的单值性定义一个 oper
局部系 $\chi$。由于它不依赖于 $\xi_a\in\CH^{\sst (2)}_a$，人们可以把这一改动归因于空间
$\CH^{\sst (2)}_a$ 自身的改动；于是该空间满足 Hecke 本征值性质。

\subsection{由 $tt^*$ 方程导出 Hecke 本征值性质}\label{Hecke-4d}

本节的目标是详细展开 Hecke 函子的 $(4d+2d)$ 描述。由于 $\Om$-形变把涨落局域化到圆作用的
固定点，人们可以把极限 $\ep_2\ra 0$ 看作把半球 $B^4_{\ep_1\ep_2}$ 退紧化为
$B^2_{\ep_1}\times\BR^2$，其中 $B^2_{\ep_1}$ 是一个二维半球。我们所关心的面算符将支撑在
$B^2_{\ep_1}$ 上。

目标是给式\eqref{NSlim+} 的行为一个四维的解释。回想：配分函数
$\mathcal{Z}_v(a;x,\tau,t;\ep_1,\ep_2)$ 可以解释为带 $\Om$-形变的规范场论在四维半球
$B^4_{\ep_1\ep_2}$ 上的配分函数，其定义用到了这样的边界条件：它把规范场路径积分在
$B^4_{\ep_1\ep_2}$ 上生成的状态投影为 Wilson 圈的本征态 \cite{NW,TV13}。变量 $a$ 参数化
Wilson 圈的本征值，而该 $4d$ 超对称规范场论的紫外规范耦合的集合记作 $\tau$。变量 $x$
参数化定义余二维面算符的边界条件 \cite{FGT}，而变量 $t$ 表示 GLSM 的 FI 参数 $t$，
该 GLSM 实现额外的余四维面算符 \cite{AGGTV}。

\subsubsection{因子化}

当 $\ep_2\ra 0$ 时，面算符上的自由度不可能对自由能的发散有贡献，因为此极限下出现的
二维椭球最终具有有限体积。因此我们预期 $\CF=\log\CZ$ 表现为
\begin{equation}
\mathcal{F}\sim-\frac{1}{\ep_2}\mathcal{Y}^{\rm 4d}(a;\tau;\ep_1)+\mathcal{W}^{\rm 2d}(a,\tau;x,t;\ep_1)+\dots\,,
\end{equation}
其中 $\mathcal{W}^{\rm 4d}(a;\tau;\ep_1)$ 与二维耦合 $x$、$t$ 无关。与式\eqref{NSlim} 比较给出
\[
\mathcal{W}^{\rm 2d}(a,\tau;x,t;\ep_1)=\mathcal{W}^{\rm\sst M2}(a,\tau;t;\ep_1)
+\mathcal{W}^{\rm\sst M5'}(a,\tau;x;\ep_1)\,,
\]
其中 $\mathcal{W}^{\rm\sst M2}(a,\tau;t;\ep_1)=\log\psi_v(t;a,\tau;\ep_1)$，
$\mathcal{W}^{\rm\sst M5'}(a,\tau;x;\ep_1)=\log\Psi(x;a,\tau;\ep_1)$。这种解耦大概可以这样
理解：定义余二维面算符的边界条件只涉及体场，而不涉及用来描述余四维面算符的自由度。

于是很自然地把 $\mathcal{Z}^{\rm 2d}_v(a,\tau;t;\ep_1):=\psi_v(t)$ 解释为：GLSM 位于下方的
二维半球 $B^2_{\ep_1}$ 上、带 $\Om$-形变、耦合到体内的 $\CS$ 类理论、并服从记作
$\mathfrak{b}_v$ 的边界条件时的配分函数。这类配分函数的例子近来在 \cite{FKNO} 中研究过，
那里发现它与 \cite{HO,HR} 所研究的半球配分函数有密切关系。人们观察到
$\mathcal{Z}^{\rm 2d}_v(a,\tau;t;\ep_1)$ 对 FI 参数 $t$ 的依赖是全纯的。此外，该配分函数与
GLSM 的拓扑扭曲版本中的一个重叠有关 \cite{HO,HR}：
\begin{equation}\label{Zoverlap}
\mathcal{Z}^{\rm 2d}_v(a,\tau;t;\ep_1)\,=\,\langle\,\mathfrak{b}_v\,|\,0\,\rangle\,,
\end{equation}
其中 $|\,0\,\rangle$ 是半球上无算符插入时的路径积分所生成的状态，而
$\langle\,\mathfrak{b}_v\,|$ 是与边界条件 $\mathfrak{b}_v$ 相联系的边界态。

\subsubsection{$tt^*$ 方程}

重叠 $\langle\,\mathfrak{b}_v\,|\,0\,\rangle$ 有一个自然的推广：把 $|\,0\,\rangle$ 换成另一个
超对称基态 $|\,1\,\rangle$，后者由手征环生成元 $v$ 从 $|\,0\,\rangle$ 得到，
$|\,1\,\rangle= v|\,0\,\rangle$。\cite{HIV} 中已指出，$\Pi_i^v:=\langle\,\mathfrak{b}_v\,|\,i\,\rangle$
表示 Cecotti 与 Vafa \cite{CV} 的 $tt^*$-联络的一个水平截面。在我们的情形中，只有一个耦合
$t$ 和一个二维的超对称真空空间。因此 $tt^*$-联络具有平坦 $SL(2,\mathbb{C})$-联络的形式
$\nabla'=(\partial_t+\mathcal{A}_t)dt$，$\nabla''=(\bar\partial_{\bar{t}}+\mathcal{A}_{\bar{t}})d{\bar{t}}$
\cite{Ga09}。由 \cite{FKNO,HO,HR} 中观察到的全纯性可得：在所讨论的情形中，
$\mathcal{A}_{\bar{t}}$ 在基 $\Pi_i^v$ 下必须幂零。这是存在一个保持 $\Pi_0^v$、并把联络
$\mathcal{A}$ 与某个 oper 相联系的规范变换的必要条件。水平性条件
$\nabla' \Pi^v=0=\nabla''\Pi^v$ 于是给出 $\Pi_0^v$ 的一个二阶微分方程。以上观察为
\cite{DFKMMN} 中关于 opers 的描述提供了一个物理实现。

由 \cite{AGGTV} 中给出的论证可得，耦合 $t$ 可以等同于 Riemann 面 $C_\tau$ 上的局部坐标。
从物理上应当预期：由水平性条件得到的这个常微分方程在 $C_\tau$ 上应当没有奇点。这将意味着
$tt^*$-联络必须规范等价于一个 oper 联络。\cite{FKNO,HO,HR} 的结果支持这一论断。

\subsubsection{固定单值性}

当 $C_\tau$ 的复结构处在 Riemann 面模空间的边界附近——在那里 $C_\tau$ 可表示为一组裤
（pairs of pants）用长管 $T_1,\dots,T_{3g-3+n}$ 连接起来——$\CS$ 类四维理论存在弱耦合
的拉格朗日描述。若 $t$ 位于某条特定的管 $T_r$ 上，则耦合的 $4d$–$2d$ 系统是弱耦合的。
该拉格朗日描述中含有 GLSM 与四维体理论之间 $ta_r/\ep_1$ 形式的耦合，其中 $a_r$ 是与 $T_r$
相联系的多重态中标量场在余四维面算符支撑上的限制 \cite{AGGTV}。因此，与 $t\ra t+2\pi$
相应的配分函数 $\mathcal{Z}^{\rm 2d}_v(a;\tau;t;\ep_1)$ 的单值性是对角的，并可由乘以
$e^{2\pi va_r/\ep_1}$ 表示。这固定了 $tt^*$-联络单值性的一半；而由于该联络规范等价于一个
oper 联络，这一半信息已足以把它完全固定下来。

于是 $tt^*$-联络被等同于几何 Langlands 对应一侧出现的 oper 局部系 $\rho$。把 ${\rm Bun}_G$
上整体微分算子代数对由 $D_r-E_r$（$r=1,\dots,3g-3+n$）生成的理想取商，便定义了 Langlands
对应另一侧的相应 $\CD$-模 $\De_{\rho}$ \cite{Fr07}。因子化式\eqref{NSlim+} 可以看作 Hecke
本征值性质在与 $\CD$-模 $\De_{\rho}$ 相联系的微分方程之解的层面上的体现。

\bigskip
\noindent\rule{\textwidth}{0.6pt}
% [Tutor Note: 临界水平的比例约定（易错）。关系式 $k=-h^{\vee}-\ep_2/\ep_1$ 必须与第 5 节的 $b^2=\hbar$ 以及第 7 节的 $\hbar=\ep_1/\ep_2$ 这一套约定配套理解：于是 $k=-h^{\vee}-1/b^2$，$\ep_2\ra0$ 即 $\hbar\ra\infty$、$b\ra\infty$，$k\ra-h^{\vee}$，确实是临界水平。若有人取相反约定 $b^2=\ep_2/\ep_1$，就会把临界水平错配到 $\ep_1\ra0$ 的极限，并连带把式\eqref{NSlim} 指数中的 $1/\ep_2$ 误写成 $1/\ep_1$。另须注意：临界水平的发散因子是 $\frac{1}{k+h^{\vee}}=-\hbar=-\ep_1/\ep_2$，因此 $\mathcal{Y}$ 必须按 $\mathcal{Y}=\ep_1\cdot(\text{经典 Yang--Yang 函数})+\CO(\ep_1^2)$ 归一化，才能使 $\mathcal{Y}/\ep_2$ 在同时标度 $\ep_i\ra\lambda\ep_i$ 下不变；核对办法是取 $\ep_1\ra0$ 应回到 \cite{NRS} 的经典 Yang--Yang 函数。]
% [Tutor Note: $E_r=\pa_{\tau_r}\mathcal{Y}_\si$ 极易被误读为 $E_r=\pa_{a_r}\mathcal{Y}_\si$。两点提示：(i) 本节后文明确以 $\tau$ 记“该 $4d$ 超对称规范场论的紫外规范耦合的集合”，故 $\tau_r$ 是第 $r$ 个裤分解分量的耦合，而 $\pa_a\mathcal{Y}$ 给出的是与长度坐标 $a$ 共轭的复 Fenchel--Nielsen 坐标（即上文 $(a,\pa_a\mathcal{Y})$ 中的第二个分量）；(ii) 该式是经典 Matone 关系 $u=\pa\mathcal{F}/\pa\tau$ 在 NS 极限下的类比，其成立依赖 $\mathcal{Y}$ 的归一化（即式\eqref{NSlim} 中 $1/\ep_2$ 的约定）。若误取 $E_r=\pa_{a_r}\mathcal{Y}$，则 $u=u^\sharp+\sum_r E_r\vartheta_r$ 不再是 Hitchin 基上的坐标，式\eqref{NSlim} 的“几何 Langlands 具体体现”一说也就随之失效。]
% [Tutor Note: 原文在 $tt^*$ 一节跳过了关键一步。$\mathcal{A}_{\bar{t}}$ 在基 $\Pi_i^v$ 中幂零，正是“存在保持 $\Pi_0^v$ 的规范变换把 $tt^*$-联络化为 oper 形式”的条件：oper 联络的 $(0,1)$ 部分必须是严格上三角的幂零矩阵（在 $\Pi_0^v$ 为最高权向量的规范下），这样水平性条件给出的二阶方程才不含 $\bar\partial_{\bar t}$ 项、成为 $C_\tau$ 上真正的二阶常微分方程 $(\ep_1^2\pa_t^2-u)\Pi_0^v=0$。而“该方程在 $C_\tau$ 上无奇点”这一步，原文只以“从物理上应当预期”带过；数学上它对应 \cite{DFKMMN} 中 oper 与 Hitchin 方程（$SL(2)$ 自对偶方程在共形极限下）之间的等价，需另行验证，不能视为已证。]
% [Tutor Note: 单值性因子 $e^{2\pi va_r/\ep_1}$ 需核对三处。(1) 没有 $i$：$t$ 是复化（全纯）的 FI 参数，耦合 $ta_r/\ep_1$ 对 $t$ 全纯，故 $t\ra t+2\pi$ 给出实指数的乘性因子，而非 $e^{2\pi i(\cdots)}$；若把 $t$ 写成 $t=\theta+i\xi$ 型实/虚分解，这个 $2\pi$ 对应 $\theta$ 方向的周期。(2) 分母是 $\ep_1$ 而非 $\ep_2$：出现在耦合里的 $\ep_1$ 来自沿面算符支撑方向的 $\Om$-形变参数，与式\eqref{NSlim} 指数中的 $\ep_2$ 是两个不同的参数，不可互换。(3) 原文把解空间基的指标 $v=\pm$ 当作整数荷 $v=\pm1$ 写进指数，这是记号混用（见本小节前文“指标 $v=+,-$ 标记解空间的一组基”）；若要严格，应说明该荷来自 $W_{1/2\hbar}$ 插入的权。此外“固定一半单值性即足以完全固定 $tt^*$-联络”隐含了 $SL(2,\mathbb{C})$ 平坦联络由单值表示（模共轭）决定这一 Riemann--Hilbert 事实，并须先假设该 oper 联络确由这些单值性实现。]
% [Tutor Note: 本节原文有三处笔误/记号不一致，译文一律保留原样，请读者注意。(a) 式(6.3) 写的是 $\mathcal{Y}^{\rm 4d}$，紧接着的一句却称“$\mathcal{W}^{\rm 4d}(a;\tau;\ep_1)$ 与二维耦合无关”——按上下文该处应为 $\mathcal{Y}^{\rm 4d}$，$\mathcal{W}$ 与 $\mathcal{Y}$ 混用了。(b) 式\eqref{NSlim} 中写作 $\Psi(a;x,\tau;\ep_1)$，而式\eqref{NSlim+} 中写作 $\Psi(x;a,\tau;\ep_1)$，自变量次序不一致（正文又说二者“就是”同一函数）。(c) 式\eqref{NSlim+} 左端漏掉了下标 $v$，而右端 $\psi_v$ 带 $v$。另，“面算符上的自由度不可能贡献自由能发散，因为二维椭球体积有限”这一论证只是启发式的：$1/\ep_2$ 极点实际来自四维体理论在 $\Om$-背景下的局域化（每个瞬子贡献在 $\ep_2=0$ 处的极点），而面算符部分给出的是 $\CO(\ep_2^0)$ 的正规项。]
